\section{Projective detuning and weak-interaction singularity}

\subsection{Proof of Theorem \ref{thm:phasecapacity}}

For $D=d+1$, scalar-extend the balanced character lattice to the
$d$-dimensional rational space
\[
 P_D^{\Q}=\left\{q\in\Q^D:\sum_jq_j=0\right\}.
\]
The detuning map $\mathcal D_E:P_D^{\Q}\to\R$ is rational-linear.  Its
kernel is the rational span of the balanced integer relation lattice.  Its
range is the rational span of anchored gaps: each generator $E_j-E_1$ is the
image of the balanced character $e_j-e_1$, while every balanced rational
combination reduces to a combination of these anchored gaps.  Rank--nullity
therefore gives
\[
 \dim_{\Q}\operatorname{range}\mathcal D_E+
 \dim_{\Q}\ker\mathcal D_E=D-1.
\]

Pontryagin duality identifies characters fixed by the completed phase law
with integral characters in the detuning kernel.  The operational temporal
dimension is consequently the range rank.  In a replica expansion, a
balanced character of positive mass $h$ first appears at order $h$; phase
cycling annihilates every character of larger mass.  This gives the final
height-filtration statement.

\subsection{Proof of Theorem \ref{thm:interaction}}

Write the product index as
$\{0,\ldots,m\}\times\{0,\ldots,n\}$.  The uncoupled energy is
$E_{ab}=A_a+B_b$.  Its gap space is generated by the $m$ anchored gaps of
$A$ and the $n$ anchored gaps of $B$, so its temporal rank is at most $m+n$.

Write a diagonal rank-one interaction as
$V_{ab}=\alpha_a\beta_b$.  For each nonzero balanced integer character $k$,
its perturbed detuning is a real polynomial in $(\alpha,\beta)$.  This
polynomial is not identically zero: if its uncoupled value vanishes, choose
one pair $(a,b)$ with $k_{ab}\ne0$ and switch on only $\alpha_a$ and
$\beta_b$.  The resulting quadratic coefficient is $k_{ab}$.  The complement
of its zero set is open and dense.  Intersecting these complements over the
countable family of nonzero balanced characters leaves a dense set in every
ball about the origin.  Choose $(\alpha,\beta)$ there with
$\max_{a,b}|\alpha_a\beta_b|\le\delta$.  Then $V=A\otimes B$ is diagonal,
$\|V\|\le\delta$, and the perturbed spectrum is maximally projective.  Its
temporal rank is $(m+1)(n+1)-1=mn+m+n$, proving the rank jump.

For self-adjoint $H,H'$ and $U(t)=e^{-itH}$,
$U'(t)=e^{-itH'}$, Duhamel's formula gives
\[
 U'(t)-U(t)=-i\int_0^t e^{-i(t-s)H'}(H'-H)e^{-isH}\,ds
\]
with reversed orientation for $t<0$, and hence
$\|U'(t)-U(t)\|\le |t|\|H'-H\|$.  Splitting
$U'(t)^*OU'(t)-U(t)^*OU(t)$ into two terms and using unitary invariance gives
$2|t|\|H'-H\|\|O\|$, hence \eqref{eq:duhamel-main}.  Finally, expansion of the $r$th fidelity moment
shows that its non-permutation excess is a sum over unequal multiplicity
vectors of equal total energy.  Maximal projective rank removes every such
balanced equality, leaving exactly the permutation floor.

\section{Arbitrary interaction graphs}

\subsection{Generic component rank}

For a connected two-site graph, the allowed edge term is the full Hermitian
matrix space on the tensor product.  Choose a simple eigenbasis for which
every one-site transition graph is connected and non-bipartite, and choose a
maximally projective eigenvalue list.  The failures of these conditions are
proper real-algebraic sets, so such choices form a nonempty open set inside
the simple-spectrum locus.

Attach a leaf to a tree.  Start from a product extension with layer shifts
chosen to avoid energy collisions.  A small edge perturbation joins the
layer transition graphs.  For every old observer site, the first-order
response of a cross-layer matrix element is a rational function of a layer
gap.  One selected transition contributes a nonzero simple-pole residue;
maximal projective rank prevents cancellation by another oriented gap.
Avoiding the finitely many zeros of these response functions makes every old
transition graph connected while preserving an odd cycle.  The local
matrix algebra at the new leaf supplies its connectivity and odd cycle.

Now suppose a nonconstant balanced rational eigenvalue relation vanished
identically on the enlarged tree-local space.  The second-order response to
the leaf interaction has simple poles indexed by oriented edges of the old
transition graph.  Vanishing residues force relation coefficients to agree
along the categorical product of that graph with the leaf complete graph.
Connectedness and non-bipartiteness make this product connected, so every
coefficient is equal, contradicting nonconstancy.  Thus no balanced relation
polynomial vanishes identically.

Symmetrizing over spectral labels turns each candidate relation into a
polynomial in local Hamiltonian coefficients.  Removing its proper zero set,
the discriminant, and the countable family of rational relations gives a
dense residual full-measure maximal set on every tree.  A spanning tree
embeds into every connected graph, so the same conclusion holds for any
connected graph.

For a disconnected graph, component gap modules add.  Independently scaling
component witnesses avoids the countable set of ratios at which two modules
intersect.  This proves the component formula
\eqref{eq:graphrank-main}.  There are only finitely many induced subsystems;
intersecting their pullbacks supplies one simultaneous generic set.

\subsection{Curvature, density and exhaustion}

On one vertex, temporal rank is $q_v-1$.  On two vertices it is
$q_uq_v-1$ when the edge is present and
$(q_u-1)+(q_v-1)$ otherwise.  Subtraction gives
\eqref{eq:curvatureedge}.  For a partition, refine every connected component
into its connected fragments of dimensions $d_1,\ldots,d_s$.  Its curvature
contribution is
\[
 \prod_{j=1}^s d_j-1-\sum_{j=1}^s(d_j-1),
\]
which is nonnegative for $d_j\ge2$ and vanishes exactly when $s=1$.  Summing
over components proves Theorem \ref{thm:curvature}.

For connected $G$, the anchored phase vector is Haar-equidistributed on
$\T^{D-1}$.  A chordal coordinate lies within $\varepsilon$ of one precisely
on an arc of normalized length
$2\arcsin(\varepsilon/2)/\pi$.  Product Haar measure proves
\eqref{eq:exactdensity}.  Substituting $D=q^n$ and taking logarithms gives
\eqref{eq:thermograph} and \eqref{eq:doubleloggraph}.

For an infinite coefficient product and a finite-volume exhaustion, pull
back each finite-volume generic set along coefficient restriction.  These
maps are continuous open surjections.  The countable intersection of the
residual pullbacks is residual and dense.  Absolute continuity of every
finite marginal makes each bad cylinder null; countable subadditivity gives
the simultaneous probability-one statement.  Since $q_v\ge2$ and volumes
grow strictly, the product dimension minus one tends to infinity.

\section{Replica collisions and component laws}

For positive weights $w_i$, expand the $r$th survival-fidelity moment into
pairs of multiplicity vectors $p,q$ with $|p|=|q|=r$.  A term survives
completed-time averaging exactly when
\[
 \sum_i(p_i-q_i)E_i=0.
\]
Maximal projective rank forces $p=q$, which is precisely the permutation
floor.  Combining this equivalence with the path genericity theorem proves
Theorem \ref{thm:pathcollision}.

For a tensor-sum spectrum, every total energy is a sum of component energies.
Mutual transversality separates a surviving equality component by component;
maximality inside each component then forces equality of the two empirical
histograms.  This proves \eqref{eq:histcollision}.  For uniform weights, the
number of $m$-tuples with histogram $k$ is $\binom{m}{k}$, giving
\eqref{eq:adm} and \eqref{eq:sffproduct}.  At fixed $m$, distinct-label
tuples contribute $m!+O(d^{-1})$, while repeated-label tuples vanish in the
limit.  The factorial moments determine the unit-mean exponential law;
independence of component Haar coordinates gives the product law.

\section{Choi targets and static complexity}

\subsection{Projective Choi tube}

The full $D$-torus splits measure-preservingly into one common phase and a
zero-anchored $(D-1)$-torus.  Equal-weight Choi densities are invariant under
the common phase, so the projective target has exactly the measure of its
full-torus lift.

If a Hilbert--Schmidt target ball is nonempty, choose one captured phase
point.  Every other captured point lies within twice the radius of it.
Pure-state density geometry converts this into a lower bound on normalized
phase overlap.  After rotating the common phase, the inequality
$1-\cos\theta\ge2\theta^2/\pi^2$ puts the relative angles inside a Euclidean
ball.  The standard ball-volume and Gamma bounds give $B_D(\delta)$.  A union
over at most $M^{s+1}$ padded words proves the Haar bound.  Maximal projective
rank transfers it to upper completed-time occupation.  Theorem
\ref{thm:path} supplies that rank on one dense residual full-measure set for
every connected interaction graph, proving Theorem \ref{thm:choicap}.

\subsection{Pauli alternation and hard polygons}

For bounded-weight Pauli coefficients outside the projective relation zero
sets, the Choi energy list is maximally projective.  The strict inequality
in \eqref{eq:paulibudgets} leaves a nonempty open complement to all padded
circuit targets.  Minimality of the projective flow makes both the identity
neighbourhood and this hard complement syndetic.  Applying them successively
gives return--hardness--return within one uniform window length.

The local separator uses the Pauli transfer matrix.  The Hamiltonian
coefficient vector is fixed by autonomous evolution, so the corresponding
local transfer block has Frobenius norm at least one.  Product-stabilizer
twirling weights every output Pauli by $3^{-\operatorname{wt}}$.  Hence the
two-query signal is at least $3^{-q}/M_q$, which exceeds $1/(D+1)$ under the
second inequality of \eqref{eq:paulibudgets}.

For Theorem \ref{thm:primepolygon}, sample a phase vector on the $p$-torsion
grid.  When $2r<p$, congruence of replica multiplicities modulo $p$ is
ordinary equality, so the $2r$th target-overlap moment is at most
$r!D^{-r}$.  Markov's inequality and a union bound over
$(p-1)|Y|$ target--vertex pairs leave one phase vector with every nonzero
multiple hard.  The $p$th multiple is identity.  A nonempty open
neighbourhood of this polygon is visited syndetically by every maximally
projective phase flow.

\section{Countable completion and reflection}

For a countable maximally projective spectrum, every finite restriction has
maximal projective rank.  Its annihilator is therefore trivial, so the orbit
closure is the full product torus.  Removing any bounded initial time does
not change this closure.  Two independent coordinates rule out a nonzero
exact global return.  For a positive summable phase-defect probe, uniform
tail control reduces approximation to finitely many coordinates.  Compact
minimality then gives a bounded hitting window.  Connectedness of the image
and continuity of the probe turn approximate level visits into exact visits
for every interior level.

Absolutely summable diagonal perturbations can be chosen arbitrarily small
and maximally projective on every finite restriction.  They are trace class,
preserve the maximal diagonal domain, and preserve self-adjointness.  A
separable perturbation $u_i v_j$ has the same properties on a product
carrier.  Along its affine coupling ray, each nonzero balanced character can
vanish at at most one coupling unless it annihilates both bare spectrum and
direction.  The chosen direction has no such nonzero character, leaving only
a countable exceptional set.  This proves Theorem \ref{thm:countable}.

Finally, autonomous composition gives
$U_{R-t}=U_RU_t^*$.  Inverting and reversing a word preserves its length in
an inverse-closed dictionary.  Choi-metric unitary invariance and the
triangle inequality add only the return defect $\Delta_R$, proving both
inequalities in \eqref{eq:reflection}.  When $\Delta_R=0$, the two bounds are
inverse inequalities and give exact palindromic symmetry.
